\documentclass[../../main.tex]{subfiles}
\graphicspath{{\subfix{../../image/}}}

\begin{document}

\section{输运方程}\label{section:ch02-transport}

最简单的偏微分方程之一是常系数输运方程。它是如下 PDE
\begin{align}\label{align:ch02-transport}
u_t+b\cdot Du=0 \qquad \text{在 }\mathbb{R}^n\times(0,\infty)\text{ 中},
\end{align}
其中 $b$ 是 $\mathbb{R}^n$ 中取定的向量，$b=(b_1,\cdots,b_n)$，$u:\mathbb{R}^n\times[0,\infty)\to\mathbb{R}$ 是未知函数，$u=u(x,t)$。这里 $x=(x_1,\cdots,x_n)\in\mathbb{R}^n$ 表示空间中一个典型点，$t\geqslant 0$ 表示一个典型时刻。我们记 $Du=D_xu=(u_{x_1},\cdots,u_{x_n})$ 为 $u$ 关于空间变量 $x$ 的梯度。

哪些函数 $u$ 是\eqref{align:ch02-transport} 的解呢？为回答此问题，暂且假定已给某个光滑解 $u$，并设法把它算出来。为此，我们首先必须认识到：偏微分方程\eqref{align:ch02-transport} 断言 $u$ 的某个方向导数为零。我们利用这一洞察，任取一点 $(x,t)\in\mathbb{R}^n\times(0,\infty)$，定义
\begin{align*}
z(s)\triangleq u(x+sb,t+s) \qquad (s\in\mathbb{R}).
\end{align*}
于是计算得
\begin{align*}
\dot{z}(s)=Du(x+sb,t+s)\cdot b+u_t(x+sb,t+s)=0 \qquad \left(\dot{\ }=\frac{\mathrm{d}}{\mathrm{d}s}\right),
\end{align*}
其中第二个等号由\eqref{align:ch02-transport} 得到。因此 $z(\cdot)$ 是 $s$ 的常值函数；从而对每个点 $(x,t)$，$u$ 在过 $(x,t)$、方向为 $(b,1)\in\mathbb{R}^{n+1}$ 的直线上为常数。于是，若我们知道 $u$ 在每条这种直线上任一点的值，就知道它在 $\mathbb{R}^n\times(0,\infty)$ 中处处的值。

\subsection{初值问题}\label{subsection:ch02-transport-ivp}

因此，为确定起见，考虑初值问题
\begin{align}\label{align:ch02-transport-ivp}
\begin{cases}
u_t+b\cdot Du=0 & \text{在 }\mathbb{R}^n\times(0,\infty)\text{ 中},\\
u=g & \text{在 }\mathbb{R}^n\times\{t=0\}\text{ 上}.
\end{cases}
\end{align}
这里 $b\in\mathbb{R}^n$ 与 $g:\mathbb{R}^n\to\mathbb{R}$ 已知，问题是求出 $u$。

给定如上 $(x,t)$，过 $(x,t)$、方向为 $(b,1)$ 的直线可参数化为 $(x+sb,t+s)$（$s\in\mathbb{R}$）。当 $s=-t$ 时，此直线到达平面 $\Gamma\triangleq\mathbb{R}^n\times\{t=0\}$，交点为 $(x-tb,0)$。由于 $u$ 在该直线上为常数，且 $u(x-tb,0)=g(x-tb)$，我们推出
\begin{align}\label{align:ch02-transport-solution}
u(x,t)=g(x-tb) \qquad (x\in\mathbb{R}^n,\ t\geqslant 0).
\end{align}
因此，若\eqref{align:ch02-transport-ivp} 有一个足够正则的解 $u$，它必定由\eqref{align:ch02-transport-solution} 给出。反之，直接可验证：若 $g$ 是 $C^1$ 的，则由\eqref{align:ch02-transport-solution} 定义的 $u$ 确实是\eqref{align:ch02-transport-ivp} 的解。

\textbf{弱解.} 若 $g$ 不是 $C^1$ 的，则\eqref{align:ch02-transport-ivp} 显然没有 $C^1$ 解。但即便在这种情形下，公式\eqref{align:ch02-transport-solution} 也肯定提供了一个很强的、事实上也是唯一合理的候选解。因此，即使 $g$ 不是 $C^1$ 的，我们也可以非正式地宣布 $u(x,t)=g(x-tb)$（$x\in\mathbb{R}^n$，$t\geqslant 0$）为\eqref{align:ch02-transport-ivp} 的一个弱解。即便 $g$ 以及 $u$ 不连续，这一切仍然成立。非光滑甚至不连续的函数有时也可以是某个 PDE 的解，这一概念在后文 \S3.4 研究非线性输运现象时还会再次出现。

\subsection{非齐次问题}\label{subsection:ch02-transport-nonhom}

接下来考察相应的非齐次问题
\begin{align}\label{align:ch02-transport-nonhom}
\begin{cases}
u_t+b\cdot Du=f & \text{在 }\mathbb{R}^n\times(0,\infty)\text{ 中},\\
u=g & \text{在 }\mathbb{R}^n\times\{t=0\}\text{ 上}.
\end{cases}
\end{align}
如前取定 $(x,t)\in\mathbb{R}^n\times(0,\infty)$，受上述计算的启发，令 $z(s)\triangleq u(x+sb,t+s)$（$s\in\mathbb{R}$）。则
\begin{align*}
\dot{z}(s)=Du(x+sb,t+s)\cdot b+u_t(x+sb,t+s)=f(x+sb,t+s).
\end{align*}
因而
\begin{align*}
u(x,t)-g(x-tb)=z(0)-z(-t)=\int_{-t}^{0}\dot{z}(s)\,\mathrm{d}s
=\int_{-t}^{0}f(x+sb,t+s)\,\mathrm{d}s
=\int_0^t f(x+(s-t)b,s)\,\mathrm{d}s,
\end{align*}
于是
\begin{align}\label{align:ch02-transport-nonhom-sol}
u(x,t)=g(x-tb)+\int_0^t f(x+(s-t)b,s)\,\mathrm{d}s \qquad (x\in\mathbb{R}^n,\ t\geqslant 0)
\end{align}
是初值问题\eqref{align:ch02-transport-nonhom} 的解。

后文 \S2.4.1 求解一维波动方程时，我们将用到此公式。

\begin{remark}\label{remark:ch02-特征法注记}
注意到，我们实际上是通过把偏微分方程转化为常微分方程，才推导出解\eqref{align:ch02-transport-solution} 与\eqref{align:ch02-transport-nonhom-sol} 的。这一做法是特征线方法的一个特例，该方法将在 \S3.2 中展开。
\end{remark}

\end{document}